\section{A gap below full terminal-path saturation}
\label{sec:second-order-path-gap}

Paths and cycles are simple.
Parallel edges remain distinct, and a pair of parallel edges can form a cycle.
The term \emph{saturated} means that the terminal-path count equals $2^{\beta(G)}$.

\begin{lemma}\label{lem:full-two-pole-classification}
Let $N$ be a connected loopless multigraph with two distinct degree-two terminals $a,c$.
Every other vertex has degree three.
Suppose that $P_{ac}(N)=2^{\beta(N)}$.
Then its bridges join saturated bridgeless two-terminal pieces in series.
Every such piece is either a pair of parallel terminal edges, or a capped permutation ladder.
The ladder permutation is a concatenation of singleton blocks and adjacent transpositions.
\end{lemma}

\begin{proof}
Every terminal path meets every cycle.
This follows by counting its affine parity class, which has $2^{\beta(N)}$ members.
Each member is one terminal path together with vertex-disjoint cycles.

Every shore of a bridge contains a cycle.
If a shore has $n$ vertices and $p$ original terminals, its rank is $(n-p+1)/2\ge1$.
Hence every bridge separates $a$ from $c$.
Neither original terminal meets a bridge.
A path from such a terminal would immediately cross it and avoid a cycle on its starting shore.
All bridges occur on their route, and no cyclic part can lie off that route.
After deleting the bridges, the attachment vertices of each piece are distinct.
Otherwise, every terminal path would avoid a cycle inside that piece.
The pieces are two-terminal graphs with degree-two attachment vertices and all other degrees three.
Their ranks add, and their terminal-path counts multiply.
Saturation forces every piece to be saturated.

Consider one bridgeless piece.
A connected bridgeless graph of maximum degree three has no cut vertex.
Choose two internally vertex-disjoint terminal paths $P,Q$.
Let $F$ contain the unused edges.
The terminals have degree zero in $F$.
Other vertices of $P\cup Q$ have degree one in $F$, while unused vertices have degree three.
The graph $F$ has no cycle, since such a cycle would avoid $P$.
No forest component can have two leaves on the interior of $P$.
Its joining path would form a cycle with a subpath of $P$, disjoint from $Q$.
The same argument applies to $Q$.
Thus every nonempty component of $F$ is one edge joining the two path interiors.

The two path interiors therefore have the same size $k$, and together they contain all nonterminal vertices.
Their unused edges give a perfect matching.
When $k=0$, the piece consists of two parallel terminal edges.
Otherwise, order both interiors from $a$ to $c$ and write the matching as a permutation $\pi$.
A pattern $231$ or $321$ produces a two-rung cycle on its first two positions.
The third rung joins the top right interval to the bottom left interval outside that cycle.
Together with $a,c$, these intervals give a terminal path disjoint from the cycle.
A pattern $312$ gives the symmetric contradiction using its last two positions for the cycle.
Thus $\pi$ avoids these three patterns.
If its first value is at least three, the later values one and two give a $312$ or $321$ pattern.
If its first value is two, value one must occur second, since an intervening value would give $231$.
Delete the initial singleton or transposition and apply induction.
The common endpoints $a,c$ cap the two ends of the ladder.
\end{proof}

\begin{lemma}\label{lem:one-full-three-pole-gap}
Let $T$ be connected, with three distinct degree-two terminals $a,b,c$ and all other degrees three.
If $\beta(T)=h$ and $P_{ac}(T)=2^h$, then
\[
 P_{bc}(T)=2^h
 \quad\hbox{or}\quad
 P_{bc}(T)\le\frac78\,2^h.
\]
The factor $7/8$ is attained by a simple graph.
\end{lemma}

\begin{proof}
Suppress $b$ to an edge of a saturated two-terminal graph $N$.
Its two neighbors are distinct.
Otherwise, its parallel-edge cycle would be pendant and would avoid every $a$--$c$ path.
Apply Lemma~\ref{lem:full-two-pole-classification}.

Suppose first that $b$ subdivides a bridge.
The earlier pieces have total positive rank, while every $b$--$c$ path avoids them.
The normalized count is therefore at most $1/2$.
The same bound holds if $b$ lies in any piece after the first.
In that case, the later pieces multiply both the local path count and its upper bound by the same power of two.
The earlier pieces contribute at least one additional rank unit only to the denominator.
It remains to consider an edge in the first piece.
A piece consisting of two parallel edges gives a triangle after subdivision, and its local count is saturated.

Now use the capped ladder description of the first piece.
Suppose that it has $k$ ladder positions, so its rank is $k+1$.
For a prefix containing $p$ positions and the left cap $a$, let $A_p$ count paths between its two right attachment vertices.
Set $A_0=1$ for the empty prefix.
For $p\ge1$, this prefix has rank $p$, so $A_p\le2^p$.
A suffix containing $q$ positions and the right cap has $2^q$ paths from either left attachment to its terminal.
Any path from $b$ to the right terminal uses exactly one edge of the suffix's two-edge entrance cut.
Its crossing parity is odd, and a simple path cannot use either edge twice.
Thus the suffix contributes one terminal path, and the following local lists are exhaustive.

First suppose that the subdivided edge is an interface edge or an end-cap edge.
With $p$ positions before it and $q$ after it, the path count is
\[
 (1+A_p)2^q,
 \qquad
 \frac{P_{bc}}{2^{k+1}}=\frac{1+A_p}{2^{p+1}}.
\]
The path first goes right, or first uses a prefix cap and returns along the other interface edge.
The displayed ratio is one when $p=0$, and at most $3/4$ when $p\ge1$.

If the subdivided edge is the rung of a single-rung block, its two initial directions give
$2(1+A_p)2^q$ paths.
Here $k=p+1+q$, so the same normalized ratio applies.

Finally, suppose that the subdivided edge lies in a square block.
Delete that edge, and regard a path through the earlier prefix as one virtual cap between the square's left ports.
There are four local paths to the right ports that avoid this cap.
There are three that use it.
For clarity, label the square ports $u,v$ on the left and $x,y$ on the right, and subdivide $ux$ by $b$.
The cap-free paths are
\[
 bx,\quad bu yv x,\quad bu y,\quad bx v y.
\]
The paths using the virtual cap $uv$ are
\[
 bu v x,\quad bu v y,\quad bx v u y.
\]
Each virtual cap has $A_p$ realizations.
Consequently,
\[
 P_{bc}=(4+3A_p)2^q,
 \qquad
 \frac{P_{bc}}{2^{k+1}}=\frac{4+3A_p}{2^{p+3}}.
\]
This ratio is $7/8$ when $p=0$, and at most $5/8$ when $p\ge1$.
All edge types have now been covered.

For equality, take the capped ladder with one square block and subdivide any square edge.
Its old terminal-path count is eight, and the new count is seven.
The graph is simple.
\end{proof}

\begin{theorem}\label{thm:terminal-path-saturation-gap}
Let $G$ be a connected finite graph or multigraph of rank $b$, with distinct terminals $s,t$.
Then
\[
 P_{st}(G)=2^b
 \quad\hbox{or}\quad
 P_{st}(G)\le\frac{15}{16}\,2^b.
\]
The factor $15/16$ is attained by a simple graph of rank four.
\end{theorem}

\begin{proof}
Induct on $b$.
For $b=0$, the graph is a tree and has one terminal path.
Use the block route between $s,t$.
An off-route cyclic block contributes rank but no terminal choices, so it gives $P_{st}(G)\le2^{b-1}$.
Otherwise, the route-block ranks add and their terminal-path counts multiply.
If there are at least two cyclic route blocks, each has smaller rank.
Induction proves the dichotomy by multiplying their normalized counts.
Bridge blocks contribute one and can be suppressed.
Thus it remains to consider a two-connected loopless block with distinct terminal attachments.

At any vertex $v$ of this block, the incident-edge pattern of a uniform affine terminal-parity vector is uniform subject to its required parity.
Indeed, deleting $v$ leaves a connected graph.
Every compatible incident pattern extends to a parity vector there, with the same number of extensions.
At an internal vertex of degree $d$, a simple path uses zero or two incident edges.
If $d\ge4$, this bounds its density by
\[
 \frac{1+\binom d2}{2^{d-1}}\le\frac78.
\]
At a terminal of degree $d\ge3$, only singleton incident patterns are allowed, giving density at most
$d/2^{d-1}\le3/4$.
Both cases prove the required bound immediately.

Suppress unmarked degree-two vertices.
We may therefore assume that both terminals have degree two and every other vertex has degree three.
The rank-one case is a cycle and is saturated.
Assume $b\ge2$.
If the terminals are adjacent, their direct edge and the remaining paths give
$P_{st}(G)\le1+2^{b-1}\le(3/4)2^b$.
Otherwise, let $x,y$ be the neighbors of $s$.
They are distinct.
If both edges joined $s$ to one vertex, that vertex would be a cut vertex unless the graph were the rank-one digon.

The graph $T=G-s$ is connected and has rank $b-1$.
Its three distinct terminals $x,y,t$ have degree two, and every other vertex has degree three.
Split paths according to their first edge:
\[
 P_{st}(G)=P_{xt}(T)+P_{yt}(T).
\]
Apply induction to both terminal pairs in $T$.
If neither count is full, their sum is at most $(15/16)2^b$.
If one is full, Lemma~\ref{lem:one-full-three-pole-gap} makes the other either full or at most $(7/8)2^{b-1}$.
Thus their sum is either $2^b$ or at most
$(1+7/8)2^{b-1}=(15/16)2^b$.

For equality, use the rank-three three-terminal example at the end of Lemma~\ref{lem:one-full-three-pole-gap}.
Add a new terminal $s$ joined to its two left terminals, and keep its right terminal $t$.
The two branch counts are eight and seven.
The resulting simple graph has rank four and fifteen terminal paths.
\end{proof}

In particular, a sequence of four-terminal blocks whose four cross-path densities tend to one is eventually exactly saturated.
Thus vanishing cross-path defects cannot bypass Theorem~\ref{thm:saturated-ring-barrier}.
